% Options for packages loaded elsewhere
\PassOptionsToPackage{unicode}{hyperref}
\PassOptionsToPackage{hyphens}{url}
%
\documentclass[
  11pt,
]{article}
\usepackage{lmodern}
\usepackage{amssymb,amsmath}
\usepackage{ifxetex,ifluatex}
\ifnum 0\ifxetex 1\fi\ifluatex 1\fi=0 % if pdftex
  \usepackage[T1]{fontenc}
  \usepackage[utf8]{inputenc}
  \usepackage{textcomp} % provide euro and other symbols
\else % if luatex or xetex
  \usepackage{unicode-math}
  \defaultfontfeatures{Scale=MatchLowercase}
  \defaultfontfeatures[\rmfamily]{Ligatures=TeX,Scale=1}
\fi
% Use upquote if available, for straight quotes in verbatim environments
\IfFileExists{upquote.sty}{\usepackage{upquote}}{}
\IfFileExists{microtype.sty}{% use microtype if available
  \usepackage[]{microtype}
  \UseMicrotypeSet[protrusion]{basicmath} % disable protrusion for tt fonts
}{}
\makeatletter
\@ifundefined{KOMAClassName}{% if non-KOMA class
  \IfFileExists{parskip.sty}{%
    \usepackage{parskip}
  }{% else
    \setlength{\parindent}{0pt}
    \setlength{\parskip}{6pt plus 2pt minus 1pt}}
}{% if KOMA class
  \KOMAoptions{parskip=half}}
\makeatother
\usepackage{xcolor}
\IfFileExists{xurl.sty}{\usepackage{xurl}}{} % add URL line breaks if available
\IfFileExists{bookmark.sty}{\usepackage{bookmark}}{\usepackage{hyperref}}
\hypersetup{
  hidelinks,
  pdfcreator={LaTeX via pandoc}}
\urlstyle{same} % disable monospaced font for URLs
\usepackage[margin=1in]{geometry}
\setlength{\emergencystretch}{3em} % prevent overfull lines
\providecommand{\tightlist}{%
  \setlength{\itemsep}{0pt}\setlength{\parskip}{0pt}}
\setcounter{secnumdepth}{-\maxdimen} % remove section numbering
\usepackage{esint}
\setmathfont{Asana Math}
\setlength{\emergencystretch}{3em}

\author{}
\date{}

\begin{document}

\hypertarget{marked-determinant-charge-and-signed-first-order-cost}{%
\section{Marked determinant charge and signed first-order
cost}\label{marked-determinant-charge-and-signed-first-order-cost}}

27 September 2026. Continues \texttt{fbe5404} and the I40--I49 proofs.
The external reviews are preserved under \texttt{guidance/source/};
their proposed general no-loss theorem is a research target, not a
premise here.

The new connection is between a marked \textbf{whole smooth core} and
the actual determinant connection, without requiring a nonzero spinor
inside the rescaled core. Its relative charge is controlled at every
finite tree depth. This is not yet control of the opposite-spinor second
Chern number: an explicit classical family distinguishes those charges.
A separate signed localization below identifies the flux needed to
improve the balanced small-h coefficient.

\hypertarget{a-marked-determinant-atom}{%
\subsection{1. A marked determinant
atom}\label{a-marked-determinant-atom}}

Use the complex (physical) orientation of \(\Gamma_n=\frac1n(1,-2)\),
where \(n=2k+1\ge3\). The generator acts as
\((z_1,z_2)\mapsto(\zeta z_1,\zeta^{-2}z_2)\). Put

\[
 f_n=\frac{k}{2k+1}=\frac{n-1}{2n}\in[1/3,1/2).
 \tag{1.1}
\]

Let \(\Omega\) be a compact oriented smooth four-manifold with this
single marked boundary. Let \(L\) be a Hermitian line bundle whose
boundary flat character, in this marking, is \(\zeta\) or
\(\zeta^{-1}\). No topology or metric restriction is imposed on the
interior. Write \(F_A=i\omega_A\), with \(\omega_A\) real, for a unitary
connection.

\hypertarget{the-relative-square-and-its-sign}{%
\subsubsection{1.1. The relative square and its
sign}\label{the-relative-square-and-its-sign}}

The minimal resolution has intersection matrix \(Q_k\) with diagonal
\((-2,\ldots,-2,-3)\) and adjacent off-diagonal entries 1. Its
anticanonical class evaluates as \((0,\ldots,0,-1)\); its boundary
character is \(\zeta^{-1}\). Inverting the last column gives

\[
 (Q_k^{-1})_{ik}=-\frac{i}{2k+1},\qquad
 (Q_k^{-1})_{kk}=-f_n.
 \tag{1.2}
\]

Indeed the interior recurrence is \(x_{i-1}-2x_i+x_{i+1}=0\), with
\(x_0=0\) and \(x_{k-1}-3x_k=1\). The inverse line has the same square.
Glue an arbitrary same-marked filling to the reversed resolution,
matching the flat boundary connections. The resulting closed line bundle
has integral first-Chern square. Consequently, if \(A\) is flat on a
boundary collar,

\[
 q_A:=\frac1{4\pi^2}\int_\Omega\omega_A\wedge\omega_A
                  \in\mathbb Z-f_n.
 \tag{1.3}
\]

The statement uses the actual character, not just that it is faithful. A
different character can have a different fractional square.

For a general boundary connection let \(A_0\) be the flat connection on
the same boundary line. Write \(A=A_0+ia\) there and set

\[
 b_A(Y)=\frac1{4\pi^2}\int_Y a\wedge da,
 \qquad q_A^{\rm rel}=q_A-b_A(Y).
 \tag{1.4}
\]

These quantities are unambiguous: \(H^1(Y;\mathbb R)=0\), so changing
the flat reference or gauge changes \(a\) by an exact form and does not
change the integral. Flatten the connection near the boundary, or use
the abelian transgression \(d(a\wedge da)=da\wedge da\), to obtain

\[
 \boxed{q_A^{\rm rel}\in\mathbb Z-f_n,\qquad
 |q_A^{\rm rel}|\ge f_n.}
 \tag{1.5}
\]

Boundary orientation is outward throughout. In particular, the
flat-collar template satisfies

\[
 \boxed{\int_\Omega|F_A|^2\ge4\pi^2 f_n
             =2\pi^2\frac{n-1}{n}\ge\frac{4\pi^2}{3}.}
 \tag{1.6}
\]

This is a characteristic-number energy bound, not an SW existence
theorem. An actual positive-weight SW solution with identically flat
determinant connection on an open collar would have zero spinor there.
The application therefore uses the general-boundary estimate below, not
a flat actual collar.

\hypertarget{arbitrary-actual-connection-on-a-controlled-geometric-collar}{%
\subsubsection{1.2. Arbitrary actual connection on a controlled
geometric
collar}\label{arbitrary-actual-connection-on-a-controlled-geometric-collar}}

Suppose an outer collar is identified with
\([r,2r]\times S^3/\Gamma_n\), with its metric uniformly bilipschitz to
the Euclidean annulus and its radial coarea factors uniformly
comparable. The comparison constant is fixed; \(r\), \(n\), interior
geometry and internal scale separation are unrestricted.

For every smooth level \(Y_s\), take the round Coulomb primitive of the
restricted curvature. On the unit round quotient, the relevant positive
Hodge spectrum is bounded below by that of \(S^3\): lift to invariant
forms, solve on \(S^3\), and descend. Thus the constant is independent
of \(n\). Scaling and metric comparison give

\[
 \|a\|_{L^2(Y_s)}\le Cr\|F_A|_{Y_s}\|_{L^2(Y_s)},\qquad
 |b_A(Y_s)|\le Cr\int_{Y_s}|F_A|^2.
 \tag{1.7}
\]

Coarea supplies a level with

\[
 |b_A(Y_s)|\le C\int_{\text{collar}}|F_A|^2.
 \tag{1.8}
\]

Let \(\Omega_s\) be the whole region inside that level and let
\(\widehat\Omega\) include it and the collar. Combining (1.4), (1.8) and
\(|\omega_A\wedge\omega_A|\le|F_A|^2d\mu\) proves

\[
 \boxed{f_n\le |q_A^{\rm rel}(\Omega_s)|
               \le C\int_{\widehat\Omega}|F_A|^2.}
 \tag{1.9}
\]

The relative number is independent of the selected level. No smallness
of the connection, no gauge equation and no field-amplitude premise is
required for (1.9). The geometric collar, marked character and whole
filling are essential hypotheses. A terminal punctured region does not
meet them.

\hypertarget{sharpness-of-the-flat-comparison-template}{%
\subsubsection{1.3. Sharpness of the flat comparison
template}\label{sharpness-of-the-flat-comparison-template}}

\href{https://arxiv.org/html/math/0206229}{Calderbank--Singer, Theorem
A} supplies scalar-flat Kähler ALE metrics on the minimal resolution.
Their canonical determinant curvature is anti-self-dual in the physical
orientation, has relative square \(-f_n\), and hence has energy exactly
\(4\pi^2f_n\). Truncate sufficiently far out and flatten the decaying
connection to its flat reference over the outer annulus. The extra
energy tends to zero. If desired, flatten the metric collar as in I47.
Homothety does not change the four-dimensional curvature-square energy.
Therefore

\[
 \inf_{\text{flat-collared marked smooth fillings}}
                   \int|F_A|^2=4\pi^2 f_n.
 \tag{1.10}
\]

This invokes a classical construction; it is not a new ALE metric or a
closed actual zero-end escape sequence.

\hypertarget{actual-determinant-charge-is-tight-on-small-volume-marked-cores}{%
\subsection{2. Actual determinant charge is tight on small-volume marked
cores}\label{actual-determinant-charge-is-tight-on-small-volume-marked-cores}}

Retain the actual finite-pencil input and class-square hypothesis of
\href{20260927_conformal_separation_and_zero_end.md}{I41, Sections
6--7}. For \(\beta\ge4\) and \(0<\eta\le1\), write

\[
 \eta=\kappa_\beta+\varepsilon,\qquad
 \epsilon_*=\beta^{-1}+\sqrt\eta.
\]

The positive determinant fit and the negative determinant saturation
give bounded two-form models \(K_+,K_-\) with

\[
 \|K_+\|_\infty+\|K_-\|_\infty\le C,\qquad
 \int(|F_A^+-K_+|^2+|F_A^--K_-|^2)\le C\epsilon_*V.
 \tag{2.1}
\]

Equivalently for the present purpose, there is a positive measure
\(\nu\) such that

\[
 |F_A|^2d\mu\le C\,d\mu+\nu,\qquad
                        \nu(M)\le C\epsilon_*V.
 \tag{2.2}
\]

This follows pointwise from the squared distance to the bounded models.
It is normalized-volume tightness, not local scale-invariant regularity.

Let \(\widehat\Omega_i\) be any finite collection of the marked
whole-core regions of Section 1.2 with overlap at most \(\Lambda\). Put
\(Q_i=q_A^{\rm rel}(\Omega_i)\). Summing (1.9) and (2.2) yields

\[
 \boxed{\frac1V\sum_i|Q_i|
 \le C\frac{\sum_i\operatorname{Vol}(\widehat\Omega_i)}V
                          +C\Lambda\epsilon_*,
 \qquad \frac N3\le\sum_i|Q_i|.}
 \tag{2.3}
\]

In a zero-end sequence, bounded overlap and total core volume \(o(V)\)
therefore force the \textbf{absolute relative determinant-charge
density}, not just the number density, to tend to zero. No nonzero
rescaled spinor in a core is assumed. The actual line and its boundary
marking are retained.

Alternatively suppose each enlarged core has volume at most a
sufficiently small fixed \(v_*\). The volume term in the unnormalized
(2.3) is at most \(CNv_*\), and can be absorbed using
\(N/3\le\sum_i|Q_i|\). Thus

\[
                        N/V\le C\Lambda\epsilon_*.
 \tag{2.4}
\]

For \textbf{absolute} exclusion of every such core, one needs the
corresponding absolute error tending to zero; for example, total core
volume tending to zero and \(\epsilon_*V\to0\). Normalized estimates
alone do not give that.

\hypertarget{what-survives-an-arbitrarily-deep-finite-decomposition}{%
\subsubsection{2.1. What survives an arbitrarily deep finite
decomposition}\label{what-survives-an-arbitrarily-deep-finite-decomposition}}

For a region with one outer and several inner marked boundaries define
the following quantity. Each inner boundary is oriented outward from the
removed \textbf{child core}, opposite to its induced boundary
orientation in the region:

\[
 Q_A(W)=\frac1{4\pi^2}\int_W\omega_A\wedge\omega_A
          -b_A(Y_{\rm out})+\sum b_A(Y_{\rm in}).
 \tag{2.5}
\]

Matched restrictions of the \textbf{same actual connection} give exact
additivity when pieces are glued. Internal boundary terms cancel before
absolute values are taken. A one-child flat neck can have zero charge
and zero energy; do not charge it as if it were a smooth cap.

Apply (2.3) to a disjoint frontier of whole subtrees, not to every
nested vertex. Its constant is independent of tree depth. A sequence of
finite trees of unbounded depth therefore obeys the same frontier
estimate. This is a proved depth-independent accounting result
\textbf{given} those marked regions. It does not construct a bubble tree
from weak integral input, nor license an infinite unsigned sum over
overlapping ancestors.

More explicitly, set \(\phi_e=b_e-f_{n_e}\) at every marked interface. A
shell with outgoing edge \(e_0\) and incoming edges \(e_1,\ldots,e_s\)
satisfies

\[
 q_A(W)=m_W+\phi_{e_0}-\sum_{a=1}^s\phi_{e_a},\qquad m_W\in\mathbb Z.
 \tag{2.6}
\]

Sum this identity over any finite retained tree. Only its outer root and
unresolved frontier remain. Equations (1.8), (2.2), and (2.3) bound the
\textbf{absolute} sum of frontier \(b_e\), \(f_{n_e}\), and hence
\(\phi_e\) by \(o(V)\) when the frontier whole cores have total volume
\(o(V)\). They also control integer charges: if \(Q_i=m_i-f_{n_i}\),
then \(|m_i|\le2|Q_i|\). This supplies the limiting estimate beyond
formal finite-tree cancellation.

\hypertarget{the-same-relative-charge-in-the-actual-harmonic-class}{%
\subsubsection{2.2. The same relative charge in the actual harmonic
class}\label{the-same-relative-charge-in-the-actual-harmonic-class}}

Let \(\mathcal H_A=2\pi i\,a_g^{\rm h}\) be harmonic curvature on the
\textbf{same integral determinant line}, including its torsion
restriction, and let \(J_+=\|F_A^+\|_2^2-4\pi^2p\). The supplied
actual-field estimates give \(J_+\le C\eta V\). Hodge orthogonality and
Stokes imply

\[
 \boxed{\|F_A^\pm-\mathcal H_A^\pm\|_2^2=J_+,\qquad
              \|F_A-\mathcal H_A\|_2^2=2J_+.}
 \tag{2.7}
\]

Indeed the difference is exact, say \(i\,d\alpha\), and
\(\int d\alpha\wedge d\alpha=0\). Its two chiral energies agree.
Combining (2.1) and (2.7) gives the same measure majorant and setwise
tightness for \(\mathcal H_A\) as for \(F_A\).

If two connections on this line differ by \(i\alpha\), the change of
\(\int\omega_A\wedge\omega_A\) is the boundary integral of
\(2\alpha\wedge\omega_A+\alpha\wedge d\alpha\). This is precisely the
change of the boundary expression in (1.4). Therefore

\[
            q_{A_{\rm h}}^{\rm rel}(\Omega_i)=q_A^{\rm rel}(\Omega_i).
 \tag{2.8}
\]

There is no estimate for the global primitive \(\alpha\), nor any use of
a first eigenvalue bound. The equality is topological; the energy
control comes from (2.7). Replacing the integral line by only its real
cohomology vector would lose the torsion boundary marking and is not
permitted.

\hypertarget{density-in-the-original-budget}{%
\subsubsection{2.3. Density in the original
budget}\label{density-in-the-original-budget}}

Now assume \(p\ge A>0\) and \(\chi(M)/A\le B\). At small Ricci residual
\(r=\fint|\operatorname{Ric}+g|^2\), Chern--Gauss--Bonnet yields

\[
 \frac{8\pi^2\chi}V\ge
 \frac{(4-2\sqrt r)^2}{24}-\frac r2
                 =\frac23-\frac23\sqrt r-\frac r3.
 \tag{2.9}
\]

Here \(\fint(s+4)^2\le4r\) and
\(\fint|\operatorname{Ric}^{\circ}|^2\le r\). For \(r\le1/4\), (2.9)
gives \(V/A\le32\pi^2B\). Positive determinant saturation gives
\(p/V\to2/(9\pi^2)\) at the endpoint. Consequently

\[
 \frac{9\pi^2}{2}(1+o(1))\le\frac VA
                         \le12\pi^2B(1+o(1)).
 \tag{2.10}
\]

In particular (2.3) is also an \(o(A)\) charge statement. The same Euler
identity bounds the total \(W^-\) energy by \(C(B)V\), hence the
curvature energy of \(E_-=S^-\otimes L\) by \(C(B)V\). None of this
gives local noncollapse or a quartic derivative bound for the torsion
field.

\hypertarget{the-determinant-charge-does-not-control-the-torsion-index}{%
\subsection{3. The determinant charge does not control the torsion
index}\label{the-determinant-charge-does-not-control-the-torsion-index}}

The distinction can be made exact. Now additionally assume a Spin-c
structure on the entire filling with determinant \(L\), extending the
specified boundary lift. An arbitrary line bundle from Section 1 need
not be such a determinant. In an exactly flat-conical comparison collar
with the two allowed parallel boundary spinors, let \(J_+\) be the
relative second Chern number of \(S^+\), and \(J_-\) that of
\(S^-\otimes L\), using these boundary sections. Write

\[
 q_A^{\rm rel}=m-f_n,\quad
 B_n=\frac{5-n}{2}+f_n,\quad
 j=2\chi(\Omega)+3\tau(\Omega)-\frac{5-n}{2}.
 \tag{3.1}
\]

Both \(m,j\) are integers. The spinor characteristic identities,
together with the outward cone Gauss--Bonnet boundary term \(1/n\), give

\[
 \boxed{J_+=\frac{m-j}{4},\qquad
        J_--J_+=2q_A^{\rm rel}+\chi(\Omega)-\frac1n
                    =\chi(\Omega)-1+2m.}
 \tag{3.2}
\]

For the first formula, \(4c_2(S^+)=c_1(L)^2-(p_1+2e)\) and the
integrated last parenthesis is \(2\chi+3\tau-B_n=j-f_n\). For the
second, the exact characteristic-form identity is
\(c_2(S^-\otimes L)-c_2(S^+)=2c_1(L)^2+e\). The parallel boundary
references have zero transgression. These derivations also specify the
sign and the necessary congruence \(m\equiv j\pmod4\).

On the canonical minimal-resolution models used in (1.10),

\[
 \chi=k+1,\quad\tau=-k,\quad m=j=J_+=0,
 \qquad J_-=k=\frac{n-1}{2},\qquad
 |q_A^{\rm rel}|=f_n<\frac12.
 \tag{3.3}
\]

Thus even a zero positive-spinor relative index and uniformly bounded
determinant energy do not bound the opposite-spinor index. The latter
grows without bound. Actual nonparallel boundary sections can
additionally carry integer homotopy defects; they must be tracked rather
than silently replaced by the flat references.

Consequently (2.3) cannot alone eliminate the order-V signed torsion
flux forced by the actual zero-end identities. One needs actual control
of that index or flux, a uniform charge-per-core bound, or a different
direct PDE argument. The classical model proves this limitation of a
proposed \textbf{charge comparison}, not existence of a closed zero-end
counterexample.

\hypertarget{an-actual-relative-index-escape-criterion}{%
\subsubsection{3.1. An actual relative-index escape
criterion}\label{an-actual-relative-index-escape-criterion}}

The existing zero-end identities give

\[
 \frac{c_2(E_-)[M]}V
   =\frac1{8\pi^2V}\int|W^-|^2
                      +\frac5{12\pi^2}+O(\epsilon_*).
 \tag{3.4}
\]

Suppose \textbf{disjoint} marked whole cores contain all zeros of the
actual cubic torsion \(t_\Phi\), which is nonzero on their boundaries.
Let \(\nu_i\) be the relative second Chern index defined by this actual
boundary section, not by a reference parallel spinor. Excision then
gives \(\sum_i\nu_i=c_2(E_-)[M]\). Bounded overlap alone would not
justify this index identity.

If their total volume is \(o(V)\), (2.3) implies, for every fixed
\(L_0\),

\[
 \frac1V\sum_{\nu_i^+\le L_0|Q_i|}\nu_i^+\longrightarrow0.
 \tag{3.5}
\]

Thus macroscopic positive torsion index must occur at unbounded
\(\nu_i^+/|Q_i|\), or outside this proposed marked-core cover. Equation
(3.3) shows why topology alone does not bound that ratio.

There is also an \textbf{actual local gauge} reason not to impose a
universal index/charge comparison on arbitrary cores. In the existing
hyperbolic ball construction of
\href{20260927_conformal_separation_and_zero_end.md}{I41, Section 15.1},
the determinant line is trivial, so Stokes gives exactly
\(q_A^{\rm rel}=0\). The actual cubic torsion has relative index \(+1\),
and \(\int|F_A|^2=O(\varepsilon^4)\). Ricci and Weyl errors vanish on
that ball. Its boundary winding and large outgoing derivative still have
to be matched to an exterior. The normalized full form defect on the
shrinking ball was not shown to tend to zero. This is a local
obstruction to a charge-only argument, not a closed O10 counterexample.

\hypertarget{actual-hyperbolic-cores-with-unbounded-torsion-index}{%
\subsubsection{3.2. Actual hyperbolic cores with unbounded torsion
index}\label{actual-hyperbolic-cores-with-unbounded-torsion-index}}

The preceding local obstruction extends beyond a single unit of charge.
It can carry arbitrarily large index on the same ball topology, with a
uniformly nonzero torsion field on a \textbf{round fixed-ratio collar}.

\textbf{Proposition (high-index local actual fields).} Fix a hyperbolic
metric with sectional curvature \(-1/3\) and a point. For every integer
\(d\ge1\) and sufficiently small \(\varepsilon>0\), there is a smooth
actual solution on the physical ball \(B_{2\varepsilon}\), on its
trivial determinant line, such that

\[
 D_A\Phi=0,\qquad F_A^+=2i\sigma(\Phi),\qquad
 |\Phi|^2=\frac43+O_d(\varepsilon^2).
 \tag{3.6}
\]

On every round sphere of radius \(s\varepsilon\), \(1\le s\le3/2\), its
cubic torsion is nonzero and has relative index \(d^2\). The relative
determinant square is zero. Uniformly over these radii,

\[
 \int_{Y_{s\varepsilon}}G=-d^2+O_d(\varepsilon^4),
 \qquad q_A^{\rm rel}=0.
 \tag{3.7}
\]

Here \(G\) is the fixed torsion-cutoff transgression of Section 5, which
equals the full unit-section transgression throughout that collar.

\textbf{Construction.} Use the Clifford frame and curved right inverses
of I41 Section 15.1. In rescaled coordinates \(x=\varepsilon y\), set

\[
 a^2=4/3,\quad c_d=2^{(d-1)/2},\quad
 g_d(z)=\frac{c_d}{2(d+1)}
              (\bar z_1^{d+1}-\bar z_2^{d+1}),\quad
 \Phi_0=a e_1+\varepsilon g_d e_2.
 \tag{3.8}
\]

The antiholomorphic coefficient has \(D_0^+(0,g_d)=0\). Direct
substitution in the cubic torsion formula gives, in physical derivative
normalization,

\[
 t_0(\varepsilon y)=a^2c_d(\bar z_1^d,\bar z_2^d).
 \tag{3.9}
\]

Indeed the components are \(a^2(g_1+ig_2,-g_3-ig_4)\); the terms
involving \(\varepsilon g_d\) cancel by the same antiholomorphic Dirac
equations.

For each fixed \(d\), the pulled-back rescaled metric is Euclidean plus
\(O(\varepsilon^2)\) in every fixed differentiability norm. The
Dirichlet right inverses \(R_D,R_A\) used in the earlier construction
are uniformly bounded for small \(\varepsilon\). First correct the seed
by \(-R_DD_\varepsilon^+\Phi_0=O_d(\varepsilon^2)\). Then apply the
contraction

\[
 A(\Phi)=R_A(0,2\sqrt2\,i\varepsilon^2\sigma(\Phi)),\qquad
 \Phi=\widehat\Phi_0-\tfrac12R_D[c(A(\Phi))\Phi].
 \tag{3.10}
\]

On a fixed \(C^{k,\alpha}\) neighborhood, \(k\ge3\), its Lipschitz
constant is \(O_d(\varepsilon^2)\). It gives a smooth exact solution
with \(A_y=O_d(\varepsilon^2)\) and \(\Phi-\Phi_0=O_d(\varepsilon^2)\),
and hence

\[
 t(\varepsilon y)=a^2c_d(\bar z_1^d,\bar z_2^d)+O_d(\varepsilon)
                                                \quad\hbox{in }C^1.
 \tag{3.11}
\]

Constants may depend on \(d\). Boundary values are outputs of the
construction; they have not been prescribed as gluing data.

Since \(c_d^2(|z_1|^{2d}+|z_2|^{2d})\ge|y|^{2d}\), the actual torsion
has \(|t|\ge a^2/2\) on \(1\le|y|\le3/2\), for sufficiently small
\(\varepsilon\). Its squared norm is at least \(4/9\), exceeding the
cutoff threshold \(m_\beta/8\le16/81\). Thus \(\theta=1\) on that whole
collar, for every \(\beta\ge4\).

The map \((\bar z_1^d,\bar z_2^d)\) has real degree \(+d^2\): the
holomorphic power map has degree \(d^2\), and conjugation on two complex
coordinates preserves real orientation. Boundary nonvanishing and the
small perturbation preserve this degree on every indicated sphere.
Interior zeros may split or be nontransverse; no unique-zero assertion
is needed for the relative index. Curvature of \(E_-\) is \(O_d(1)\) in
physical norms, so
\(\int_{B_{s\varepsilon}}c_2(E_-)=O_d(\varepsilon^4)\). The relative
Chern formula proves (3.7). The determinant equality in (3.7) is exact
Stokes on a trivial line, independently of any norm estimate.

The construction and exact torsion equation also give

\[
 \begin{split}
 \int(|\nabla_A^2\Phi|^2+|\nabla t|^2)&=O_d(\varepsilon^2),\\
 \int(|F_A|^2+|\nabla_A\Phi|^4+|t|^4)&=O_d(\varepsilon^4),\\
 \int(|\langle\Phi,\nabla_A\Phi\rangle|^2+|D_{3A}^-t|^2)
                                             &=O_d(\varepsilon^6).
 \end{split}
 \tag{3.12}
\]

For the last estimate, the current is \(O_d(\varepsilon)\), the spinor
derivative is \(O_d(1)\), and the actual torsion Dirac identity cancels
the scalar/model terms; \(W^+=0\) here. Ricci and Weyl errors vanish
exactly because the physical metric is the fixed hyperbolic metric.

Choose \(\varepsilon_d\) diagonally, after fixing \(d\), so all costs in
(3.12) tend to zero and the boundary degree and cutoff estimates hold.
Then \(\nu=d^2\to\infty\), while the displayed absolute subcritical
costs vanish. There is no claim of a uniform-in-d implicit-function
radius.

The missing critical direction cost is visible explicitly. For the unit
seed \(u_d=t_0/|t_0|\), polar coordinates give

\[
 \frac{d^2}{r^2}\le|du_d|^2\le\frac{4d^2}{r^2}
                   \quad(1\le r\le3/2).
 \tag{3.13}
\]

The phase contribution is \(d^2(x^{d-1}+y^{d-1})/(r^2(x^d+y^d))\), where
\(x=\cos^2\vartheta\), \(y=\sin^2\vartheta\). It lies between
\(d^2/r^2\) and \(2d^2/r^2\). The remaining angular derivative is
bounded by \(d^2/r^2\): its ratio to this bound is
\((xy)^{d-1}/(x^d+y^d)^2\le1\). For example write \(x/y=z^2\); the last
inequality follows from \((z^{d-1}-1)(z^{d+1}-1)\ge0\), using \(x+y=1\).
The phase upper bound follows from \((x-y)(x^{d-1}-y^{d-1})\ge0\).
Endpoint values follow by continuity. The displayed upper bound is
deliberately relaxed. For each fixed \(d\), (3.11) transfers these
estimates to the actual unit field after decreasing \(\varepsilon\). The
diagonal can therefore retain

\[
 \int_{\varepsilon_d\le r\le3\varepsilon_d/2}|\nabla u|^4
                    \asymp d^4\longrightarrow\infty,
 \quad
 \int_{\varepsilon_d\le r\le3\varepsilon_d/2}|\nabla u|^2
                    \asymp\varepsilon_d^2d^2\longrightarrow0.
 \tag{3.14}
\]

This disproves a local signed-flux bound based only on the displayed
subcritical budgets, even with actual gauge equations and a nonzero
fixed-ratio torsion collar. It does \textbf{not} supply a fully
calibrated endpoint collar: torsion length has large radial/angular
variation, and the ASD determinant fit to the required exterior fails.
There is no closed conformal supplier, global class-square budget, or
exterior matching. It is not an O10 counterexample, nor a sharpness
example for the frontier exponent in (5.5).

\hypertarget{where-a-smaller-balanced-first-order-coefficient-must-be-paid}{%
\subsection{4. Where a smaller balanced first-order coefficient must be
paid}\label{where-a-smaller-balanced-first-order-coefficient-must-be-paid}}

This section concerns the established \textbf{balanced surgery family},
not the zero-end weighted solution. Let

\[
 M=X\#k\overline{\mathbb {CP}}^{\,2}
       \#\ell(S^1\times S^3),\quad A=c_1^2(X),\quad
 h=(k+4\ell)/A\to0,\quad U=1+h+O(h^2).
 \tag{4.1}
\]

Use a scalar-minus-four metric and one actual \textbf{ordinary},
unweighted SW solution \((A_{\rm conn},\Phi)\) supplied by the specified
class with \(a^2=A-k\) and \(p=(a_g^+)^2\ge A\). The symbol \(A\)
without a subscript continues to denote the topological budget, not the
connection. In the formulas below \(F_{A_{\rm conn}}\) is abbreviated to
\(F_A\). All fields and remainders in this section belong to this same
solution.

Put

\[
 \rho=|\Phi|^2\le4,\quad P=\int|\operatorname{Ric}^{\circ}|^2,
 \quad R=P+4V,\quad \mathcal E=4P/R,
 \quad J=\|F_A^+-(F_A^+)_H\|_2^2.
\]

The actual input and Hodge identities give

\[
 c_2(S^+)[M]=\ell,\qquad
 B:=\int|F_A^-|^2=4\pi^2(p-A+k)+J.
 \tag{4.2}
\]

Choose a fixed smooth increasing \(v(\rho)\) equal to zero for
\(\rho\le1/2\), one for \(\rho\ge1\), with \(|v'|\le4\), and put
\(\chi=(1-v)^2\). Write \(T_u\) for the exact rank-two transgression on
\(u=\Phi/\sqrt\rho\), as in I40 Section 1, and define

\[
 q_\chi=\int d\chi\wedge T_u,\quad
 P_\chi=\int\chi|\operatorname{Ric}^{\circ}|^2,\quad P_o=P-P_\chi,
 \quad B_\chi=\int\chi|F_A^-|^2,\quad B_o=B-B_\chi,
\] \[
 M_\chi=\int\chi\left(2|W^+|^2+\frac23-\frac{\rho^2}{8}\right)
                         \ge\frac{13}{24}\int\chi\ge0.
 \tag{4.3}
\]

The signed flux \(q_\chi\) is a different object from the relative
determinant square \(q_A^{\rm rel}\). In particular, the low-amplitude
region has not been identified as a marked whole surgery core.

\hypertarget{exact-spatial-balances}{%
\subsubsection{4.1. Exact spatial
balances}\label{exact-spatial-balances}}

The transgression identity \(dT_u=c_2(\mathbb F)\) and the extension of
\((1-\chi)T_u\) by zero near all spinor zeros give

\[
 \ell=\int\chi c_2(\mathbb F)+q_\chi.
\]

There is no transversality or regular-level-set assumption. For the
actual positive-spinor connection the signed Chern density is

\[
 16\pi^2c_2(\mathbb F)=
 \left(|F_A^+|^2-|F_A^-|^2-2|W^+|^2
             -\frac{s^2}{24}+\frac{|\operatorname{Ric}^{\circ}|^2}{2}\right)d\mu.
\]

Use \(s=-4\) and \(|F_A^+|^2=\rho^2/8\). This proves the two exact
identities

\[
 \boxed{P_\chi=32\pi^2\ell+2B_\chi+2M_\chi-32\pi^2q_\chi,}
 \tag{4.4}
\] \[
 \boxed{P-8\pi^2Ah=P_o+8\pi^2(p-A)+2J-2B_o
                              +2M_\chi-32\pi^2q_\chi.}
 \tag{4.5}
\]

For (4.5), substitute (4.2) in (4.4) and use \(Ah=k+4\ell\). The only
unfavorable terms are exterior ASD determinant energy and positive
signed transition flux. This specifies a spatial source of the earlier
signed remainder; it does not discard the Weyl payment tested by X06.

\hypertarget{first-order-necessary-conditions}{%
\subsubsection{4.2. First-order necessary
conditions}\label{first-order-necessary-conditions}}

Suppose \(\mathcal E=O(h)\). With \(e=R/(8\pi^2A)\), the two classical
energy walls give

\[
 1+h\le e\le\frac{U}{1-\mathcal E}=1+O(h).
\]

Thus on any subsequence on which \(\mathcal E/h\to c\),
\(P/(Ah)\to2\pi^2c\). Moreover \(U=1+15k/A-12\ell/A\) implies

\[
 \frac{k}{Ah}\to\frac29,\qquad
 \frac{\ell}{Ah}\to\frac7{36}.
 \tag{4.6}
\]

Dropping only the displayed nonnegative terms in (4.4) gives

\[
 \boxed{\liminf\frac{q_\chi}{Ah}
                         \ge\frac7{36}-\frac c{16}.}
 \tag{4.7}
\]

In particular, a coefficient below \(28/9\) requires positive signed
flux of first-order size. Conversely, the \textbf{additional} condition
\(\limsup q_\chi/(Ah)\le0\) forces \(c\ge28/9\), not four. The class
supply does not establish that sign condition.

For every \(c<4\), (4.5) instead gives

\[
 \boxed{\liminf\frac{B_o+16\pi^2q_\chi}{Ah}
                                      \ge\pi^2(4-c).}
 \tag{4.8}
\]

Thus proving both \(B_o=o(Ah)\) and \((q_\chi)_+=o(Ah)\) would suffice
for the coefficient-four lower bound. These are substantive missing
estimates, not consequences asserted here of near-minimality.

\hypertarget{the-flux-demands-an-actual-critical-mixed-current}{%
\subsubsection{4.3. The flux demands an actual critical mixed
current}\label{the-flux-demands-an-actual-critical-mixed-current}}

Set \(T=\{1/2<\rho<1\}\), \(\zeta=\langle\Phi,\nabla_A\Phi\rangle\), and

\[
 M_{\rm mix}=\int_T|\zeta|^2|\nabla_A\Phi|^2.
\]

For this ordinary solution, the exact nonnegative remainder in
\href{20260927_full_core_cost_and_extremal_metrics.md}{I47, Section 4.1}
reads

\[
 D_2/18=F_{\rm fit}+\Gamma/2+J+H,
\] \[
 \Gamma=\int\left(\frac{|d\rho|^2}{4}+|\operatorname{Im}\zeta|^2\right),
 \quad H=\frac23\int(1-x)^2(2+4x-x^2),\quad x=\rho/4,
\] \[
 F_{\rm fit}=\int|W^+-W_0|^2,\quad
 W_0=\tfrac12(|\sigma(\Phi)|^2I/3-\sigma(\Phi)\otimes\sigma(\Phi)).
 \tag{4.9}
\]

In particular \(Q:=\int(4-\rho)^2\le2D_2/3\) and \(\Gamma\le D_2/9\). On
\(T\), volume is at most \(Q/9\), the model \(W_0\) is bounded, and
\(F_{\rm fit}\le D_2/18\). If \(\Omega\) denotes trace-free
positive-spinor curvature, then \(G^2=|\Omega|^2+|F_A^+|^2\) satisfies

\[
 \int_TG^2\le C(P+D_2).
\]

The exact transgression in a unitary \((u,v)\) frame is

\[
 4\pi^2T_u=-a_{\rm conn}\wedge\mathbb F_{22}
 +\tfrac12(b\wedge\mathbb F_{12}-\bar b\wedge\mathbb F_{21})
 +a_{\rm conn}\wedge\bar b\wedge b.
\]

On \(T\), \(|d\rho|\le2|\zeta|\), \(|a_{\rm conn}|\le2|\zeta|\), and
\(|b|\le\sqrt2|\nabla_A\Phi|\). Hence Cauchy--Schwarz, without a Sobolev
or Hessian estimate, gives

\[
 \boxed{|q_\chi|\le
 C\left\{\sqrt{(P+D_2+B)M_{\rm mix}}+M_{\rm mix}\right\}.}
 \tag{4.10}
\]

The central term is controlled using \(\int_T|\zeta|^4\le M_{\rm mix}\).
Along the above first-order sequences,

\[
 D_2=36\pi^2A\{e-1-h-2(p/A-1)\}\ge0
\]

implies \(D_2,p-A=O(Ah)\), and (4.2), (4.9) give \(B=O(Ah)\). Thus a
coefficient bounded strictly below \(28/9\) forces
\(\liminf M_{\rm mix}/(Ah)>0\). This is a necessary actual
concentration, not its exclusion.

If additionally \(D_2=o(Ah)\), then
\(\operatorname{Vol}(T),\Gamma=o(Ah)\). To control the ordinary spinor's
quadratic derivative there, test
\(4\nabla_A^*\nabla_A\Phi+(\rho-4)\Phi=0\) against \(\theta(\rho)\Phi\),
where \(\theta=1\) below 1, zero above 2, and \(\theta'\le0\) is
bounded. Integration by parts gives

\[
 4\int\theta|\nabla_A\Phi|^2
  =\int\theta(4-\rho)\rho-2\int\theta'|d\rho|^2
  \le C(Q+\Gamma)\le CD_2.
 \tag{4.11}
\]

Therefore below the stated threshold,

\[
 \operatorname{Vol}(T),\ \int_T|\nabla_A\Phi|^2=o(Ah),
 \qquad \int_T|\zeta|^2|\nabla_A\Phi|^2\ge c_*Ah.
 \tag{4.12}
\]

In particular the derivative supremum diverges, since
\(M_{\rm mix}\le\sup_T|\nabla_A\Phi|^2\,\Gamma\). No estimate for a
different weighted solution was substituted here.

\hypertarget{conditional-critical-patterns}{%
\subsubsection{4.4. Conditional critical
patterns}\label{conditional-critical-patterns}}

If \(c=28/9\) and \(\limsup q_\chi/(Ah)\le0\), equality in the limiting
nonnegative balance (4.4) forces

\[
 P_o,B_\chi,M_\chi=o(Ah),\quad q_\chi=o(Ah),\quad
 \liminf\frac{B_o}{Ah}\ge\frac{8\pi^2}{9}.
 \tag{4.13}
\]

All leading Ricci error lies in the low-amplitude region, while the
indicated ASD determinant energy lies outside. Since \(\int\chi=o(Ah)\)
but \(P_\chi/(Ah)\to56\pi^2/9\),
\(\|\operatorname{Ric}^{\circ}\|_\infty\) diverges on that region.

If instead \(c=4\), \(B_o=o(Ah)\), and \((q_\chi)_+=o(Ah)\), then (4.5)
forces

\[
 P_o,p-A,J,M_\chi,q_\chi=o(Ah),\quad
 \frac{B_\chi}{Ah}\to\frac{8\pi^2}{9},\quad
 \frac{P_\chi}{Ah}\to8\pi^2.
 \tag{4.14}
\]

These are different conditional first-order patterns. Neither is a
classification of all minimizing sequences, and neither proves that the
coefficient-four sufficient conditions hold on the classical recovery
family.

\hypertarget{a-quantitative-criterion-for-torsion-flux-at-a-marked-frontier}{%
\subsection{5. A quantitative criterion for torsion flux at a marked
frontier}\label{a-quantitative-criterion-for-torsion-flux-at-a-marked-frontier}}

Return to the weighted zero-end solution. Let \(t=t_\Phi\), \(m=|t|^2\),
and use the inherited torsion cutoff \(\chi=1-\theta(m)^2\), with
\(\theta=0\) for \(m\le m_\beta/16\) and \(\theta=1\) for
\(m\ge m_\beta/8\). Here \(m_\beta\) has positive upper and lower bounds
for \(\beta\ge4\). Define \(G=\theta^2T_t\), extended by zero through
torsion zeros. This \(\chi\) is the torsion cutoff, not the
ordinary-spinor cutoff of Section 4.

With

\[
 k_A=9(|F_A^+|^2-|F_A^-|^2)+s^2/24-|\operatorname{Ric}^{\circ}|^2/2,
\]

the exact local version of the inherited signed identity is, for every
smooth domain \(D\),

\[
 2\int_D\theta^2|W^-|^2+\int_D\theta^2 k_A
   =16\pi^2\left(\int_Dd\chi\wedge T_t+
                                      \int_{\partial D}G\right).
 \tag{5.1}
\]

This follows by integrating \(dG=-d\chi\wedge T_t+\theta^2c_2(E_-)\).
Boundary terms cancel for an actual partition, but need not disappear in
a limiting decomposition.

\hypertarget{a-scale-free-trace-bound}{%
\subsubsection{5.1. A scale-free trace
bound}\label{a-scale-free-trace-bound}}

Consider a controlled collar of radius \(r_e\) and order \(n_e\) as in
Section 1.2. Put

\[
 E_e=\int_{\mathcal C_e}|F_{E_-}|^2,\qquad
 J_e=\int_{\mathcal C_e\cap\{m>m_\beta/16\}}|\nabla u|^4,
 \qquad u=t/|t|.
\]

Here the derivative is that of the actual unit section, including its
phase, not just its image in the projective bundle. The rank-two
transgression gives

\[
 |G|\le C\theta^2(|F_{E_-}||\nabla u|+|\nabla u|^3).
\]

Use coarea to choose an interior slice, then Hölder on the collar. Since
\(\operatorname{Vol}(\mathcal C_e)\le Cr_e^4/n_e\),

\[
 \boxed{\left|\int_{Y_e}G\right|
 \le C n_e^{-1/4}
                  (E_e^{1/2}J_e^{1/4}+J_e^{3/4}).}
 \tag{5.2}
\]

The physical radius cancels. By selecting from the good sets of the two
coarea estimates, the same slice can satisfy (1.8) and (5.2). Let
\(H=\sum_e n_e^{-1}\). Hölder in the interface index gives

\[
 \boxed{\frac1V\sum_e\left|\int_{Y_e}G\right|
 \le C\left(\frac HV\right)^{1/4}
 \left[
   \left(\frac{\sum E_e}V\right)^{1/2}
   \left(\frac{\sum J_e}V\right)^{1/4}
       +\left(\frac{\sum J_e}V\right)^{3/4}
 \right].}
 \tag{5.3}
\]

For marked frontiers of Section 2, \(H/V\le N/V\to0\). If their
normalized quartic torsion-direction energy stays bounded and their
normalized curvature energy stays bounded, no normalized torsion
boundary flux is lost at the selected cuts.

Under \(p\ge A>0\) and \(\chi(M)/A\le B\), the latter curvature bound
follows from the original actual budgets, rather than an additional
\(C^2\) collar condition. In the project's curvature norms,

\[
 |F_{E_-}|^2=|W^-|^2+s^2/48
                +|\operatorname{Ric}^{\circ}|^2/4+(9/2)|F_A|^2.
\]

Gauss--Bonnet gives \(\int|W^-|^2\le8\pi^2\chi+P/2\), and Section 2.3
controls \(\chi/V\). For bounded-overlap collars, therefore,
\(\sum E_e/V\le C(B)\Lambda\).

\hypertarget{surviving-flux-forces-a-faster-critical-concentration}{%
\subsubsection{5.2. Surviving flux forces a faster critical
concentration}\label{surviving-flux-forces-a-faster-critical-concentration}}

Write \(d=H/V\), \(j_4=\sum J_e/V\), and suppose the left side of (5.3)
is at least a fixed \(b>0\), while \(\sum E_e/V\le E_0\). At least one
of its two summands must be at least \(b/2\). Hence

\[
 j_4\ge c\min\left\{\frac{b^4}{dE_0^2},
                               b^{4/3}d^{-1/3}\right\}.
 \tag{5.4}
\]

For \(d\to0\) and fixed \(b,E_0\), this implies

\[
                 \boxed{j_4\ge c(b,E_0)d^{-1/3}.}
 \tag{5.5}
\]

In the uniformly small-core-volume version (2.4), \(d\le C\epsilon_*\),
so surviving macroscopic flux forces \(j_4\ge c\epsilon_*^{-1/3}\). This
is stronger than a merely positive quartic concentration requirement.

The quantifier matters: the macroscopic flux must survive the
\textbf{coarea-selected cuts}, or every admissible cut, not just a
preselected hypersurface. Moving a cut can transfer nonclosed \(G\) into
the bulk and transition terms of (5.1). The global signed identity does
not alone force all of its flux through a chosen marked frontier.

The lower bound for \(|t|\) on this region implies
\(|\nabla u|\le C|\nabla t|\). Current actual budgets therefore bound
the normalized quadratic direction energy, but not the quartic quantity
in (5.5). Pure radial amplitude change cannot by itself supply this
boundary flux.

For completeness put \(L_2=V^{-1}\sum\int|\nabla u|^2\) over these
collar regions. All sets below are restricted to
\(\bigcup_e(\mathcal C_e\cap\{m>m_\beta/16\})\); multiplicity is
\(\sum_e\mathbf1_{\mathcal C_e\cap\{m>m_\beta/16\}}\). With
\(L_2\le C\Lambda\) and \(j_4\to\infty\), the set
\(Z=\{|\nabla u|^2\ge j_4/(2L_2)\}\), counted with this multiplicity,
satisfies

\[
 \frac{\operatorname{Vol}_{\rm mult}(Z)}V
        \le\frac{2L_2^2}{j_4},\qquad
 \frac1V\sum_e\int_{\mathcal C_e\cap Z}|\nabla u|^4\ge j_4/2.
 \tag{5.6}
\]

Thus the forced divergent quartic mass lies on vanishing relative
volume. Equations (5.5)--(5.6) in particular force
\(\sup|\nabla u|\ge c d^{-1/6}\), and hence a corresponding bound for
the physical \(|\nabla t|\) on that region.

\hypertarget{an-explicit-routing-condition-and-the-remaining-alternative}{%
\subsubsection{5.3. An explicit routing condition, and the remaining
alternative}\label{an-explicit-routing-condition-and-the-remaining-alternative}}

Here is one sufficient condition that ensures the required macroscopic
flux survives the selected cuts. Suppose disjoint marked whole cores
have total volume \(o(V)\), all the torsion low-amplitude region
\(\{\chi\ne0\}\) lies strictly inside them, and \(\chi=0\) throughout
every admissible outer collar and the exterior. This is an additional
actual-field condition.

For any coarea-selected cuts, apply (5.1) to their exterior \(D\). Its
transition term is zero. The inherited fits give
\(V^{-1}\int_M k_A=20/3+O(\epsilon_*)\) and
\(V^{-1}\int_{M\setminus D}|k_A|=o(1)\). Since the \(W^-\) term is
nonnegative, the sum of \textbf{outward core} fluxes is negative and
obeys

\[
 -\frac1V\sum_e\int_{Y_e}G
                         \ge\frac5{12\pi^2}-o(1).
 \tag{5.7}
\]

If also \(p\ge A>0\) and \(\chi(M)/A\le B\), the curvature input to
(5.5) is automatic. The resulting alternative is precise: such a routing
decomposition of an actual endpoint sequence must have

\[
 \frac1V\sum_e\int_{\mathcal C_e\cap\{m>m_\beta/16\}}|\nabla u|^4
                 \ge c(B,K)(H/V)^{-1/3}\longrightarrow\infty.
 \tag{5.8}
\]

In particular a decomposition satisfying the strictly
weaker-than-bounded condition \(j_4=o((H/V)^{-1/3})\) is impossible.
This connects a genuine global field-routing premise to a quantified
failure, rather than assuming an arbitrary boundary has macroscopic
flux.

Weak input has not yet supplied this decomposition or ruled out (5.8).
Trivial spherical interfaces are a further distinct possibility: (5.2)
holds for them with \(n=1\), but the determinant atom does not bound
their number. The actual local example in Section 3.1 explains why they
cannot be silently omitted from a proposed decomposition. These
alternatives leave the unrestricted endpoint theorem open.

\hypertarget{why-level-cuts-alone-do-not-cancel-the-signed-trace}{%
\subsubsection{5.4. Why level cuts alone do not cancel the signed
trace}\label{why-level-cuts-alone-do-not-cancel-the-signed-trace}}

The phase part of the same transgression identifies a possible false
shortcut. On \(t\ne0\), put
\(j_t=\operatorname{Im}\langle t,\nabla t\rangle\), \(m=|t|^2\), and
\(f_A=-iF_A\). In a unit \(t\)-frame the diagonal connection is
\(i j_t/m\), and the determinant of \(E_-\) has curvature \(3i f_A\).
The simplified rank-two formula gives exactly

\[
 4\pi^2T_u=-m^{-2}j_t\wedge dj_t+3m^{-1}j_t\wedge f_A
          +\tfrac12(b\wedge\mathbb F_{12}-\bar b\wedge\mathbb F_{21}).
 \tag{5.9}
\]

The explicit radial derivative cancels from the first term because
\(j_t\wedge dm\wedge j_t=0\); the phase helicity does not cancel. Thus
taking a mass-level boundary, where tangential \(dm=0\), does not force
the trace to be small. The inherited hyperbolic local actual core
already has outward \(\int T_u=-1+O(\varepsilon^4)\) on its nonzero
outer section, while its bulk characteristic integral tends to zero. No
nonpositive-flux theorem or improved signed trace bound follows from
level selection alone.

\hypertarget{consequences-for-the-continuing-research}{%
\subsection{6. Consequences for the continuing
research}\label{consequences-for-the-continuing-research}}

There are now two distinct actual connections, not one interchangeable
charge:

\begin{itemize}
\tightlist
\item
  The determinant's full relative charge is tight on existing marked
  small whole cores, including harmonic replacement and arbitrary finite
  depth.
\item
  Torsion indices can grow relative to that charge, and an actual local
  trivial-boundary core can carry torsion index with zero determinant
  charge. Any global argument must retain the actual section framing or
  directly control the resulting signed flux.
\end{itemize}

For the zero-end target, the next sufficient step is either direct
actual control of the signed transition in (5.1), or a decomposition
that proves both a genuine flux-routing statement and enough control to
contradict (5.5). Proving the first does not require completing all
possible core classifications. For balanced first-order sharpness, (4.5)
instead identifies the two actual quantities whose leading contribution
must be decided.

The present arguments do not give an unrestricted endpoint gap, a
Gamma-limit, the balanced optimal coefficient, diameter-free
realization, or the Einstein classification of the fixed Dehn-twist
manifold. Those remain separate open statements. The written proofs
above provide specific new conditional connections and identify exactly
which proposed stronger charge comparisons are false.

\end{document}
